\section{Conclusiones y Trabajo Futuro}
La presente revisión sistemática del estado del arte ha permitido evaluar de manera integral el desarrollo, madurez y desafíos actuales de las soluciones de visión computacional y aprendizaje automático aplicadas a los sistemas de transporte inteligentes y el monitoreo de conductores. A través del análisis estadístico cuantitativo del corpus de 213 publicaciones, se evidencia una consolidación científica significativa que está transformando el control y seguridad de nuestras redes viales.

Las conclusiones principales obtenidas de este análisis global son:
\begin{itemize}
    \item La hegemonía metodológica ha migrado casi en su totalidad hacia soluciones de Deep Learning debido a su alta capacidad de generalización en entornos urbanos dinámicos.
    \item El hardware embebido (Edge Computing) es fundamental para lograr implementaciones en tiempo real que garanticen la privacidad del conductor y baja latencia de procesamiento.
    \item Los desafíos actuales se centran en mitigar la sensibilidad de los modelos ante cambios climáticos y problemas de iluminación.
\end{itemize}

Como trabajo futuro, se vislumbran tres grandes áreas de desarrollo prioritarias en la literatura científica reciente:
\begin{enumerate}
    \item \textbf{Fusión de Sensores Multimodales:} La combinación de datos visuales con señales biométricas (EEG, ECG, conductancia de la piel) o datos espaciales (LiDAR, Radar) para robustecer la detección frente a oclusiones y cambios de iluminación.
    \item \textbf{Modelos de Lenguaje-Visión (Vision-Language Models) y Aprendizaje Zero-Shot:} La incorporación de arquitecturas multimedales que permitan detectar conductas atípicas de riesgo sin depender exclusivamente de clases predefinidas de entrenamiento.
    \item \textbf{Generación de Datos Sintéticos:} El uso de Redes Generativas Adversarias (GAN) y modelos de difusión para generar imágenes realistas de conducción nocturna, oclusiones y accidentes, resolviendo el problema de escasez de datos balanceados de entrenamiento.
\end{enumerate}
